\section{Часть 1 — Введение и постановка проблемы}

\begin{frame}{Почему вопрос надёжности важен}
\begin{itemize}
  \item Большинство систем проверяются тестами.
  \item Тестирование проверяет \textbf{конечный набор} входов, а пространство входов часто огромно/бесконечно.
  \item Даже хорошо протестированная система может сломаться в редких условиях.
\end{itemize}
\end{frame}

\begin{frame}{Где это критично}
Для систем повышенной ответственности ошибки недопустимы:
\begin{itemize}
  \item операционные системы;
  \item компиляторы;
  \item криптографические библиотеки;
  \item авиационные и медицинские системы.
\end{itemize}
\vspace{2mm}
Выход — \textbf{формальная верификация программ}.
\end{frame}

\begin{frame}{Формальная верификация: что это и зачем}
\textbf{Формальная верификация} — математическое доказательство того, что программа удовлетворяет заданным свойствам.
\begin{itemize}
  \item вместо «проверили на примерах» — «доказали для всех входов»;
  \item свойства формулируются как логические утверждения;
  \item результат проверяем машинно.
\end{itemize}
\vspace{2mm}
Примеры свойств:
\begin{itemize}
  \item сортировка возвращает упорядоченный массив;
  \item компилятор сохраняет семантику программы.
\end{itemize}
\end{frame}

\begin{frame}{Цена формальной верификации}
\begin{itemize}
  \item Пример: компилятор \textbf{CompCert} (C), корректность доказана формально.
  \item Известный факт: объём доказательства может быть \textbf{существенно больше} объёма кода (в статье приводится пример порядка \(\sim 3\times\)).
\end{itemize}

\vspace{2mm}
Ключевая проблема: \textbf{трудоёмкость написания доказательств}.
\end{frame}

\begin{frame}{Proof assistant и Isabelle/HOL}
\textbf{Proof assistant} — среда для формулирования утверждений и написания \textbf{машинно-проверяемых} доказательств.

\vspace{2mm}
\textbf{Isabelle/HOL}:
\begin{itemize}
  \item Isabelle — популярный proof assistant;
  \item HOL — логика высших порядков (удобна для функций, множеств, структур);
  \item доказательства пишутся как \textbf{proof scripts}.
\end{itemize}
\end{frame}

\begin{frame}[fragile]{Мини-пример proof script (Isabelle)}
\lstset{style=code, language=}
\begin{lstlisting}
lemma add_0_right: "x + (0::nat) = x"
by simp
\end{lstlisting}

\vspace{1mm}
И пример с индукцией:
\begin{lstlisting}
lemma add_comm: "(a::nat) + b = b + a"
proof (induct a)
  case 0
  then show ?case by simp
next
  case (Suc a)
  then show ?case by simp
qed
\end{lstlisting}
\end{frame}

\begin{frame}{Автоматизация доказательств: основные подходы}
\begin{itemize}
  \item \textbf{Hammers} — подбирают факты/леммы и пытаются «пробить» доказательство эвристиками.
        В Isabelle: \textbf{Sledgehammer}.
  \item \textbf{Search-based нейропруверы} — предсказывают следующий шаг и ищут доказательство по дереву состояний.
\end{itemize}

\vspace{2mm}
Минусы search-based:
\begin{itemize}
  \item сложная архитектура (постоянный цикл «модель $\leftrightarrow$ prover»);
  \item дорогой поиск;
  \item генерация по шагам ограничивает использование больших LLM.
\end{itemize}
\end{frame}

\begin{frame}{Переход к Baldur}
Авторы предлагают другой подход:
\begin{itemize}
  \item генерировать \textbf{целое} доказательство сразу (whole-proof generation);
  \item проверять доказательство в Isabelle;
  \item при ошибке — чинить доказательство (proof repair).
\end{itemize}

\vspace{2mm}
Далее — как устроен Baldur.
\end{frame}
